\section{Discussion}
\label{sec:discussion}

\subsection{Art-Specific Conditioning vs.\ Generic Control}

Our dual-stage conditioning demonstrates that art-theoretic
principles---the temporal separation of structure and tone in
ink-wash practice---can be directly encoded as conditioning
schedules in diffusion models. This contrasts with generic
ControlNet approaches that inject all conditions uniformly
throughout the denoising process, without accounting for the
artistic workflow that the conditions represent.

\subsection{Blank-Space as a Training Objective}

The blank-space loss is, to our knowledge, the first formalization
of a negative-space aesthetic principle as a differentiable
training objective. While negative space has been discussed in
computational aesthetics, it has not been enforced as a hard
constraint on generative output. Our results suggest that explicit
blank-space preservation significantly improves the perceived
authenticity of generated ink-wash videos.

\subsection{Limitations}

\begin{itemize}
    \item \textbf{Dataset size:} 30 paintings is a small dataset
    for fine-tuning; results may not generalize to all Shanshui
    sub-styles.
    \item \textbf{Annotation cost:} line maps and ink-density maps
    require semi-manual annotation, limiting scalability.
    \item \textbf{Single backbone:} we use Wan2.1-I2V; transfer to
    other I2V backbones (LTX-Video) is not yet validated.
    \item \textbf{Short video length:} experiments use 4-frame
    clips; longer videos may require additional temporal
    consistency mechanisms.
\end{itemize}

\subsection{Future Directions}

The dual-conditioning framework generalizes beyond ink-wash: oil
painting (fat-over-lean layering), watercolor (wet-on-wet vs.\
wet-on-dry), and other traditions have analogous stage-separated
workflows that could be encoded as conditioning schedules. We also
plan to explore automatic extraction of line maps and ink-density
maps from input paintings, reducing the annotation burden.
